\vspace{4pt}
{\small % 11pt
    \begin{center}
        \textbf{ABSTRAK (Times New Roman 11, spasi 1, Maksimal 250 kata)}
    \end{center}
    \noindent Abstrak memuat uraian singkat mengenai masalah dan tujuan penelitian, metode yang digunakan, dan hasil penelitian. Tekanan penulisan abstrak terutama pada hasil penelitian. Abstrak ditulis dalam bahasa Indonesia dan Bahasa Inggris. Kata kunci perlu dicantumkan untuk menggambarkan ranah masalah yang diteliti dan istilah-istilah pokok yang mendasari pelaksanaan penelitian. Kata-kata kunci dapat berupa kata tunggal atau gabungan kata. Jumlah kata-kata maksimal 5 kata. Kata-kata kunci ini diperlukan untuk komputerisasi. Pencarian judul penelitian dan abstraknya dipermudah dengan kata-kata kunci tersebut.

    \vspace{4pt}
    \noindent\textbf{Kata Kunci:} kata1, kata2, kata3, kata4, kata5 yang relevan dengan linieritas ilmu

    \vspace{10pt}
    \begin{center}
        \textit{\textbf{ABSTRACT}}
    \end{center}
    \noindent\textit{An abstract is a brief summary of a research article, thesis, review, conference proceeding or any-depth analysis of a particular subject or discipline, and is often used to help the reader quickly ascertain the paper purposes. When used, an abstract always appears at the beginning of a manuscript or typescript, acting as the point-of-entry for any given academic paper or patent application. Abstracting and indexing services for various academic disciplines are aimed at compiling a body of literature for that particular subject. Abstract length varies by discipline and publisher requirements. Abstracts are typically sectioned logically as an overview of what appears in the paper.}

    \vspace{4pt}
    \noindent\textit{\textbf{Keywords:}} \textit{word1, word2, word3, word4, word5 related to program linier}
}

\vspace{12pt}
\noindent Naskah diterima : \tglditerima, Naskah dipublikasikan : \tgldipublikasikan